\scanpage{273}
\mspage{223}
aussi bien partir de l'\unsure{ens} de flèches $\mathrm{Liss}_W$, pour définir les localiseurs $W_S^g$ comme
\[
W_S^g=W_S^{\unsure{\mathrm{Lis}|\mathrm{Liss}}_W}\quad\text{au lieu de}\quad
W_S^g=W_S^{\unsure{\mathrm{Isoloc}|\mathrm{Iso\,loc}}},
\]
ce qui est équivalent.

D'autre part, en travaillant avec les \unsure{localisés} $\unsure{\mathrm{HOT}|\mathrm{Hot}}_W^0(S)=(\mathrm{Cat}_{/S})(W_S^g)^{-1}$, le rôle de $\underline{\mathrm{Fib}}_W S$ pour la théorie de $\mathrm{HOT}^{\unsure{\mathrm{lc}|\mathrm{li}}}_W(S)$ est joué ici par $\underline{\mathrm{Prop}}_W S$. On voit donc que dans la théorie des \inserted{$W$-}Types d'homotopie relatifs non loc. constants, notés $\mathrm{HOT}_W^0$, les ens. de flèches $\mathrm{Liss}_W$ et $\mathrm{Prop}_W$ jouent un rôle en quelque sorte complémentaire : les flèches \inserted{$f:S'\to S$} de $\mathrm{Liss}_W$ sont (pratiquement) celles qui permettent directement de définir
\[
f_W^{*0}:\mathrm{HOT}_W^0(S)\unsure{-|\longrightarrow}\mathrm{HOT}_W^0(S')
\]
comme induit par $f^*:(\mathrm{Cat}_{/S},W_S^g)\longrightarrow(\mathrm{Cat}_{/S'},W_{S'}^g)$,
tandis que les flèches de $\mathrm{Liss}_W$ permettent de définir \inserted{pour tt $S$ dans $\mathrm{Cat}$} la sous-catégorie pleine
\[
\underline{\mathrm{Prop}}_W S\subset\mathrm{Cat}_{/S}
\]
donnant lieu à une équivalence de \unsure{catégorie|catégories}
\[
\left\{\begin{aligned}
(\underline{\mathrm{Prop}}_W S)(W_S^{\unsure{u|v}})^{-1}&\xrightarrow{\approx}
\underset{\displaystyle =\,\mathrm{HOT}_W^0(S)}{\underline{(\mathrm{Cat}_{/S})(W_S^g)^{-1}}}\\
\text{car}\qquad &\\[-5pt]
W_S^u=W_S^g\ \text{sur } \underline{\mathrm{Prop}}_W S&
\end{aligned}\right.
\]
(ce qui visiblement est le substitut \unsure{annoncé} des ``axiomes de la \deleted{\unsure{co}} fibre $W$-homotopique''). Cette
% archive sheet number: 273
